\documentclass{template}

\usepackage[utf8]{inputenc}
\usepackage[ngerman]{babel}
\usepackage{hyperref}

\hypersetup{
	colorlinks=true,
	urlcolor=myblue,
	linkcolor=myblue,
	citecolor=myblue,
	pdftitle={Dr. André Kaderli -- Lebenslauf},
	pdfauthor={André Kaderli}
}

\begin{document}

	\name{Dr. André Kaderli}
	\tagline{Promovierter theoretischer Physiker (ETH Zürich, HU Berlin) mit grosser Leidenschaft für Mathematik, theoretische Physik und Data Science -- von Stringtheorie und Zahlentheorie über mathematische Modellierung und Statistik bis zu Machine Learning und KI. Seit über fünf Jahren setze ich Daten- und KI-Projekte für die öffentliche Verwaltung und die Privatwirtschaft um, mit Praxis- und Führungserfahrung von prädiktiven Modellen bis zu KI-Agenten. Besonders Freude macht es mir, komplexe Themen verständlich zu vermitteln, gerade in der Lehre von Mathematik und Physik.}
	\begin{contactline}
		Obere Bahnhofstrasse 22M, 3700 Spiez\sep
		\href{mailto:andre.kaderli@gmail.com}{andre.kaderli@gmail.com}\sep
		\href{tel:+41793323931}{+41\,79 332 39 31}\sep
		\href{https://www.linkedin.com/in/andre-kaderli/}{linkedin.com/in/andre-kaderli}
	\end{contactline}

	\datedsection{Berufserfahrung}{}

	\datedsubsection{Novalytica AG, Bern}{}

	\datedsubsubsection{Principal Data \& AI Engineer und Partner}{Seit Jan. 2026}
	\begin{itemize}
		\item Fachliche Leitung von Daten- und KI-Initiativen in öffentlicher Verwaltung und Privatwirtschaft; Führung und Mentoring interdisziplinärer Teams
		\item Konzeption, Aufbau und Betrieb moderner Daten- und Analytics-Plattformen (Microsoft, Databricks, Open Source) und Data-Governance-Frameworks; Beratung von Entscheidungsträgern
	\end{itemize}

	\datedsubsubsection{Senior Data Scientist}{Jan. 2023 -- Dez. 2025}
	\begin{itemize}
		\item Leitung von End-to-End-Datenprojekten: Datenpipelines, BI-Reports, produktive ML-Modelle und intelligente Datenagenten
		\item Statistische Analysen grosser, heterogener Datensätze und mathematische Modellierung komplexer Systeme; Verantwortung für Entwicklungsstandards (CI/CD, Testing, Monitoring)
	\end{itemize}

	\datedsubsubsection{Data Scientist}{Mai 2022 -- Dez. 2022}
	\begin{itemize}
		\item Data-Analytics- und BI-Projekte (Power BI, Azure, Databricks); Entwicklung von ML-Modellen und statistischen Auswertungen
	\end{itemize}

	\datedsubsection{Eraneos AG, Bern}{}
	\datedsubsubsection{Consultant}{Sep. 2021 -- Apr. 2022}
	\begin{itemize}
		\item IT- und Datenprojekte im öffentlichen Sektor mit Fokus auf Business-Analyse und quantitative Datenauswertungen
	\end{itemize}

	\datedsubsection{Max-Planck-Institut für Gravitationsphysik, Potsdam}{}
	\datedsubsubsection{Wissenschaftlicher Mitarbeiter (Doktorand)}{Sep. 2018 -- Aug. 2021}
	\begin{itemize}
		\item Forschung in Theoretischer Physik (Quantengravitation): analytische und numerische Methoden, rekursive Algorithmen und Simulationsverfahren zur Lösung komplexer Differentialgleichungen; Veröffentlichungen und Konferenzbeiträge
	\end{itemize}

	\datedsection{Ausbildung}{}

	\datedsubsection{PhD in Physik, Humboldt-Universität zu Berlin}{Sep. 2018 -- Aug. 2021}
	\begin{itemize}
		\item Theoretische Physik (summa cum laude)
		\item Forschung an der Schnittstelle von Stringtheorie und Zahlentheorie (Veröffentlichungen: \href{https://inspirehep.net/authors/1809024}{inspirehep.net})
	\end{itemize}

	\datedsubsection{MSc in Physik, ETH Zürich}{Sep. 2016 -- Aug. 2018}
	\begin{itemize}
		\item Theoretische und Mathematische Physik (Abschlussnote: 6.0/6.0)
		\item Fokus auf Allgemeine Relativitätstheorie, Quantenfeldtheorie und Statistische Physik
	\end{itemize}

	\datedsubsection{BSc in Physik, ETH Zürich}{Sep. 2013 -- Aug. 2016}
	\begin{itemize}
		\item Physik, Mathematik und Informatik (Abschlussnote: 5.7/6.0)
	\end{itemize}

	\newpage
	\datedsection{Lehre und Betreuung}{}

	\datedsubsection{Data Academy, Novalytica}{Seit Mai 2022}
	\begin{itemize}
		\item Kurse zu Data Literacy, Programmierung (Python, R), Statistik, Datenanalyse und Visualisierung
	\end{itemize}

	\datedsubsection{Vorlesungsunterstützung, HU Berlin}{2018 -- 2021}
	\begin{itemize}
		\item Funktionentheorie und Quantenfeldtheorie; Betreuung von Studierenden und wissenschaftlichen Arbeiten in Mathematik und Physik
	\end{itemize}

	\datedsubsection{Teaching Assistant, ETH Zürich}{2014 -- 2016}
	\begin{itemize}
		\item Übungsgruppen, Besprechung von Übungsserien und Betreuung der Studierenden in Linearer Algebra am Departement Mathematik
	\end{itemize}

	\datedsubsection{Privater Nachhilfelehrer}{2009 -- 2018}
	\begin{itemize}
		\item Mathematik, Physik und Informatik
	\end{itemize}

	\datedsection{Ausgewählte fachliche und technische Kenntnisse}{}
	\kv{Mathematik:}{Lineare Algebra, Analysis, Funktionentheorie, Wahrscheinlichkeitstheorie, Statistik, Algebra, Zahlentheorie, Geometrie, Differentialgeometrie, Topologie, algebraische Topologie}
	\kv{Angewandte Mathematik:}{Numerik, Optimierung, mathematische Modellierung und Simulation, Signalverarbeitung, stochastische Prozesse, Zeitreihenanalyse}
	\kv{Physik:}{Klassische Mechanik, Elektrodynamik, Thermodynamik und Statistische Physik, Quantenmechanik, Allgemeine Relativitätstheorie, Quantenfeldtheorie, Stringtheorie, Festkörperphysik, Quantum Computing}
	\kv{Data Science, Machine Learning \& KI:}{Analyse grosser und heterogener Datensätze, ML-Modelle für Prognosen, Mustererkennung und Anomalieerkennung, KI-Agenten und LLM-Anwendungen}
	\kv{Data Engineering \& Programmierung:}{Datenpipelines und Analyseplattformen (Azure, Databricks, Microsoft Fabric, Apache Spark, Delta Lake); Python, C++, R, MATLAB, Mathematica, SQL, PySpark}

	\datedsection{Zertifizierungen}{}
	\kv{Microsoft:}{Azure Data Scientist Associate, Azure Data Engineer Associate, Power BI Data Analyst Associate, Fabric Analytics Engineer Associate, Fabric Data Engineer Associate}
	\kv{Databricks:}{Databricks Data Engineer Professional, Databricks Data Engineer Associate}
	\kv{Projektmanagement:}{Hermes Advanced}

	\datedsection{Ausgewählte persönliche Kompetenzen}{}
	\kv{Didaktische Vermittlung:}{Komplexe fachliche Inhalte verständlich und stufengerecht aufbereiten, für Studierende ebenso wie für Fach- und Führungspublikum}
	\kv{Führung und Zusammenarbeit:}{Transparente, respektvolle Führung interdisziplinärer Teams mit Fokus auf Eigenverantwortung und klare Ziele}
	\kv{Wissenschaftliche Sorgfalt:}{Hoher Anspruch an Qualität, Nachvollziehbarkeit und Dokumentation}

	\datedsection{Weitere Angaben}{}
	\kv{Sprachen:}{Deutsch (Muttersprache), Englisch (verhandlungssicher), Französisch (Grundkenntnisse)}
	\kv{Stiftungen \& Stipendien:}{Alumni der Schweizerischen Studienstiftung, ehem. Träger des Ernst-Göhner-Stipendiums}
	\kv{Freiwilliges Engagement:}{Tourenleiter beim Schweizer Alpen-Club (SAC)}
	\kv{Interessen:}{Bergsteigen, Klettern, Skifahren, Gleitschirmfliegen, Jiu-Jitsu}

\end{document}
